% Source: https://artofproblemsolving.com/community/c3h2575008_dont_get_something
% Stable identity: aops-topic-2575008-q01
\textbf{Figure 1.} The square has side $21$. Its vertex on $AC$ cuts off two small triangles at $A$ and $C$; each has a $45^\circ$ angle and a right angle, so each is an isosceles right triangle with legs $21$. Hence $AB=BC=21+21=42$ and $AC=42\sqrt2$.

\textbf{Figure 2.} Let the square be $PQRS$ with $P,Q$ on $AC$ (in the order $A,P,Q,C$), $R$ on $BC$ and $S$ on $AB$, and let its side be $s$.

\begin{asy}
usepackage("amsmath");
size(9cm);
defaultpen(fontsize(10pt));
pair B=(0,0),A=(0,42),C=(42,0);
draw(A--B--C--cycle,linewidth(0.9));draw((0,0)--(21,0)--(21,21)--(0,21)--cycle);
label("$A$",A,N);label("$B$",B,SW);label("$C$",C,SE);label("$21$",(10.5,0),S);label("$21$",(21,10.5),E);label("Figure 1",(21,-9));
pair d=(58,0);pair A2=A+d,B2=B+d,C2=C+d,P=(14,28)+d,Q=(28,14)+d,R=(14,0)+d,T=(0,14)+d;
draw(A2--B2--C2--cycle,linewidth(0.9));draw(P--Q--R--T--cycle);
label("$A$",A2,N);label("$B$",B2,SW);label("$C$",C2,SE);label("$P$",P,NE);label("$Q$",Q,NE);label("$R$",R,S);label("$S$",T,W);
label("$s$",(A2+P)/2,NE);label("$s$",(P+Q)/2,NE);label("$s$",(Q+C2)/2,NE);label("Figure 2",(21,-9)+d);
\end{asy}

Triangle $APS$ has $\angle A=45^\circ$ and $\angle APS=90^\circ$, so it is an isosceles right triangle and $AP=PS=s$. In the same way $QC=QR=s$. This is why $AC$ is divided into three equal parts:
\[AC=AP+PQ+QC=3s,\qquad s=\dfrac{42\sqrt2}{3}=14\sqrt2.\]
The area is $s^2=\left(14\sqrt2\right)^2=392$, answer $\boxed{\textbf{(B) } 392}$.
